\section{The Three-Clock Dictionary}
\label{sec:dictionary}

Table~\ref{tab:dictionary} collects the results of Sections~\ref{sec:clock1}, \ref{sec:clock2} and~\ref{sec:clock3}. Each row is a concept from nonlinear dynamics and each column is a clock. A cell names the observed behavior that the concept describes on that clock and marks its evidence level: \statusP\ when the correspondence is proved for a model of the system, \statusM\ when it has been measured in trained networks, and \statusC\ when it is a conjecture. A question mark means that we found no work that tests the concept on that clock.

\begin{table}[p]
\centering
\scriptsize
\setlength{\tabcolsep}{3pt}
\renewcommand{\arraystretch}{1.25}
\begin{tabularx}{\textwidth}{>{\raggedright\arraybackslash}p{2.1cm} >{\raggedright\arraybackslash}X >{\raggedright\arraybackslash}X >{\raggedright\arraybackslash}X}
\toprule
\textbf{Concept} & \textbf{Clock 1: layer time $\ell$} & \textbf{Clock 2: training time $t$} & \textbf{Clock 3: generation time $n$} \\
\midrule
Fixed point & Tokens collapse to a cluster; rank collapse \statusP\ \citep{geshkovski2023mathematical,dong2021attention}. Massive activations nearly input independent \statusM\ \citep{sun2024massive} & Convergence to a minimum; neural collapse at the end of training \statusM\ \citep{papyan2020prevalence} & Period-one loops (one token repeated) as the simplest case of degenerate repetition \statusM\ \citep{holtzman2019curious} \\
Linear stability, Jacobian spectrum & Dynamical isometry; layer Jacobians of trained models \statusP/\statusM\ \citep{pennington2017resurrecting,fernando2026dynamics} & Sharpness vs.\ $2/\eta$ \statusM\ \citep{cohen2021gradient} & Derivative of the feedback map factorizes through $1/T$ \statusM\ \citep{jones2026chaos} \\
Finite-time Lyapunov exponent & Ridges in input space separate classes \statusM\ \citep{storm2023finitetime}; per-layer profile over training \statusM\ (Sec.~\ref{sec:exp1}) & Sensitivity of runs to minibatch noise early in training \statusM\ \citep{frankle2019linear} & Divergence of nearby generations \statusM\ \citep{jaca2026chaos,jones2026chaos} \\
Edge of chaos, criticality & Order-chaos transition at initialization sets trainable depth \statusP\ \citep{poole2016exponential,schoenholz2016deep} & Edge of stability \statusM\ \citep{cohen2021gradient} & Critical temperature with power-law text statistics \statusM\ \citep{nakaishi2024phase} \\
Limit cycle (periodic orbit) & \na & Period-2 orbits of GD \statusP\ \citep{chen2022edge}; SGD converges to cycles \statusP\ \citep{chaudhari2017stochastic} & Repetition loops \statusM\ \citep{xu2022learning}; paraphrase 2-cycles \statusM\ \citep{wang2025unveiling}; loops converge in hidden-state space \statusM\ (Exp.~3) \\
Period doubling, route to chaos & \na & Larger step sizes give period doubling and chaos in minimal models \statusP\ \citep{zhu2022edgeofstability,kong2020stochasticity}; period-2 orbit, then irregular oscillation, in a transformer \statusM\ (Exp.~2) & Feigenbaum scaling in embedding feedback \statusM\ \citep{jones2026chaos} \\
Fractal boundary & \na & Trainable vs.\ divergent hyperparameters \statusM\ \citep{sohldickstein2024boundary}; interleaved trainable and divergent step sizes \statusM\ (Exp.~2); fractal attractors of optimizers \statusM\ \citep{tuci2026generalization} & \na \\
Basin of attraction & Class regions bounded by FTLE ridges \statusM\ \citep{storm2023finitetime} & Which minimum is reached depends on step size and batch size \statusM\ \citep{jastrzebski2017three} & Prompts that do or do not fall into loops \statusC \\
Metastability, slow manifold & Multi-cluster states last exponentially long \statusP\ \citep{geshkovski2024dynamic} & Plateaus \statusP\ \citep{saxe2013exact}; slow phase of pretraining \statusM\ \citep{zhang2026stability,wen2024warmupstabledecay} & Long coherent stretches before drift \statusC \\
Saddle and ghost & \na & Saddles dominate \statusM\ \citep{dauphin2014identifying}; abrupt learning from a ghost \statusP\ \citep{dinc2025ghost} & \na \\
Phase transition & \na & Grokking \statusP/\statusM\ \citep{rubin2023grokking,power2022grokking}; induction heads \statusM\ \citep{olsson2022incontext} & Low-$T$ repetitive vs.\ high-$T$ incoherent phase \statusM\ \citep{nakaishi2024phase,arnold2024phase}; narrow coherent window between them, also in a reasoning model \statusM\ (Exp.~3) \\
Noise-driven (Kramers) escape & \na & Escape of SGD from minima \statusP\ \citep{xie2020diffusion,mori2021powerlaw} & Escape from repetition loops, logit gap as barrier \statusP/\statusM\ (Prop.~\ref{prop:arrhenius}, Exp.~3); locking temperature \statusP\ (Prop.~\ref{prop:tlock}), horizon-dependent in practice \statusM\ (Exp.~3) \\
Synchronization, consensus & Clustering as consensus; Kuramoto form \statusP\ \citep{geshkovski2023mathematical,li2026diverse}; synchronization by noise \statusP\ \citep{agazzi2026stochastic} & \na & Collective dynamics of interacting models \statusP\ \citep{krishnamurthy2026stochastic} \\
Koopman, linearization & Local linear description of each layer \statusM\ \citep{fernando2026dynamics} & Koopman-based training and pruning \statusM\ \citep{dogra2020optimizing,redman2021operator}; lazy regime \statusP\ \citep{jacot2018neural} & \na \\
Conservation law, symmetry & Gradient-flow structure of attention \statusP\ \citep{sander2021sinkformers} & Symmetries give conserved quantities that finite steps break \statusP\ \citep{kunin2020neural} & \na \\
\bottomrule
\end{tabularx}
\caption{The Three-Clock Dictionary. Each cell gives the behavior that a dynamical-systems concept describes on one clock, with its evidence level (\statusP: proved in a model; \statusM: measured in trained networks; \statusC: conjecture) and representative references. Question marks are cells for which we found no study; they define the open problems of Section~\ref{sec:open}.}
\label{tab:dictionary}
\end{table}

\paragraph{Reading the table by rows.} Fixed points, stability and Lyapunov exponents have been used on all three clocks, although with different meanings: on clock~1 they describe a single finite trajectory, on clock~2 the long-run behavior of an optimizer, and on clock~3 a random process. Phase transitions and escape appear on clocks~2 and~3 but not on clock~1, where the number of steps is fixed and small, so there is no long-time limit in which a transition could be sharp. Even so, the metastability result of \citet{geshkovski2024dynamic} shows that finite-time transients on clock~1 can be analyzed with care.

\paragraph{Reading the table by columns.} Clock~2 is the best covered column, because optimization theory has always been a theory of iterated maps. Clock~1 is covered mostly by results at initialization and by models with simplified attention. Clock~3 is the youngest column, and most of its entries date from 2024 or later.

\paragraph{Where the analogies are weak.} First, clock~1 runs for a fixed and small number of steps (between 6 and about 100 layers), so asymptotic notions such as attractors and Lyapunov exponents must be replaced by finite-time versions. Second, clock~2 is driven by a data stream whose statistics are not stationary in general, and by learning-rate schedules, which makes it a non-autonomous system. Third, the state of clock~3 is the whole context, so its dimension grows with time and the chain never returns to an earlier state. Statements about loops and attractors on clock~3 are therefore about the projection of the state onto recent tokens. Each of these differences limits which theorems can be imported from the classical theory without change.
